\section{Edge AI}

\subsection{Khái niệm Edge AI}

Edge AI là cách triển khai tác vụ trí tuệ nhân tạo trên thiết bị đặt gần nguồn sinh
dữ liệu, thay vì gửi toàn bộ dữ liệu thô về đám mây. Thiết bị có thể là gateway,
máy tính nhúng, camera thông minh hoặc máy công nghiệp. Lợi ích chính gồm giảm
độ trễ, giảm băng thông, tăng khả năng hoạt động ngoại tuyến và hạn chế dữ liệu
nhạy cảm rời khỏi hiện trường \cite{shi2016,satya2017}.

Edge AI không loại bỏ vai trò của Cloud. Cloud vẫn phù hợp cho huấn luyện, lưu
trữ artifact, quản lý đội thiết bị và phân tích dài hạn; Edge chịu trách nhiệm cho
vòng lặp phản hồi nhanh. Do đó, giải pháp thực tế thường là một hệ thống lai trong
đó Cloud và Edge có nhiệm vụ khác nhau nhưng chia sẻ cùng vòng đời model.

\subsection{Kiến trúc Cloud--Edge}

Kiến trúc Cloud--Edge gồm ba lớp chính: Cloud Control Plane, kênh truyền thông
và Edge Plane. Control Plane lưu desired state, model version, chính sách và lịch
sử vận hành. Edge Plane chứa Edge Agent và AI Runtime, duy trì actual state và
tiếp tục suy luận khi mất kết nối. Kênh truyền thông mang lệnh, sự kiện và
telemetry; các tệp model dung lượng lớn có thể được phân phối qua object storage.

Kiến trúc cần chấp nhận kết nối không ổn định. Lệnh phải có định danh duy nhất,
thời hạn hiệu lực và khả năng thực thi lặp an toàn. Edge Agent phải lưu trạng thái
cục bộ, đồng bộ lại khi có mạng và không được dừng runtime đang khỏe chỉ vì
Control Plane tạm thời không truy cập được.

\subsection{Đặc điểm của Edge Device}

Edge Device có tài nguyên nhỏ hơn máy chủ, thường dùng bộ nhớ chia sẻ giữa CPU
và GPU, giới hạn điện năng và tản nhiệt. Phần mềm phụ thuộc chặt vào phiên bản
hệ điều hành, CUDA, cuDNN, TensorRT và driver. Cùng một model artifact có thể
không chạy được trên hai thiết bị nếu engine được tạo cho kiến trúc GPU hoặc
phiên bản runtime khác nhau.

Thiết bị biên cũng nằm ngoài trung tâm dữ liệu nên dễ gặp mất nguồn, mất mạng,
camera lỗi hoặc thay đổi môi trường. Vì vậy, khả năng quan sát, tự khởi động lại,
rollback và xác thực artifact là một phần của yêu cầu hệ thống chứ không phải chức
năng bổ sung.

\subsection{Các vấn đề khi triển khai AI trên Edge}

Năm nhóm vấn đề chính là năng lực tính toán, bộ nhớ, latency, power và thermal.
Chúng tạo thành bài toán tối ưu đa mục tiêu: giảm số phép tính nhưng không làm
accuracy giảm quá mức; giảm precision nhưng vẫn duy trì tính ổn định số; tăng FPS
nhưng không gây quá nhiệt hoặc vượt ngân sách điện.

Ngoài model inference, hệ thống còn phải tính chi phí camera capture, chuyển đổi
màu, resize, chuẩn hóa, non-maximum suppression, truyền kết quả và telemetry.
Do đó, benchmark chỉ đo thời gian TensorRT engine là cần thiết nhưng chưa đủ để
kết luận về khả năng đáp ứng thời gian thực của toàn hệ thống.